\documentclass[10pt,a4paper]{amsart}
\usepackage[margin=19mm]{geometry}
\usepackage{amsmath,amssymb,mathtools}
\usepackage[scr=rsfs,cal=euler]{mathalfa}
\usepackage{xfrac}
\usepackage[unicode,hidelinks,bookmarks=false]{hyperref}
\newcommand{\ie}{i.e.\ }
\newcommand{\cf}{cf.\ }
\newcommand{\resp}{respectively\ }
\newcommand{\Subsection}{\subsection}
\providecommand{\nobreakdash}{\nobreak\hbox{-}}
\setlength{\emergencystretch}{2em}
\multlinegap0pt
\usepackage{mathtools}
\usepackage{amssymb}
\usepackage{enumerate}
\newtheorem{lem}{Lemma}
\newtheorem{definition}{Definition}
\newtheorem{thm}{Theorem}
\providecommand{\Linf}{\mathcal{L}_{\infty}}
\providecommand{\Lneg}{\mathcal{L}_{\infty}^-}
\providecommand{\Lpos}{\mathcal{L}_{\infty}^+}
\providecommand{\Rn}{\mathbb{R}^n}
\newcommand{\dist}{\mathrm{dist}}
\providecommand{\ue}{u_{\epsilon}}
\newcommand{\wt}{\widetilde}
\title{On the comparison principle for a nonlocal infinity Laplacian}
\author{Frida Fejne}
\date{Source: arXiv:2501.04506v2 (CC BY 4.0)}
\begin{document}
\maketitle

\noindent\emph{Source and licence.} On the comparison principle for a nonlocal infinity Laplacian. Version arXiv:2501.04506v2 (CC BY 4.0), distributed by arXiv under the Creative Commons Attribution 4.0 International licence, http://creativecommons.org/licenses/by/4.0/, as recorded in the arXiv licence metadata for this exact version and shown on the abstract page https://arxiv.org/abs/2501.04506v2. Full licence notice: \texttt{licenses/arxiv-2501.04506-CC-BY-4.0.txt}.\par\smallskip

\noindent\textit{Editorial scope.} Counted unit: the article's thm thm:comparison (source0.tex lines 358--360). The article's complete original proof of it is delivered with it (source0.tex lines 364--535, one \texttt{proof} environment of 172 lines, which the article places away from the statement). No mathematics is written, repaired or re-derived: the statement and the proof are the article's own lines, in the article's own notation, and no retained line is substituted or re-flowed.  Retained uncounted as the source-local context the proof needs: lines 162--164; lines 182--185; lines 226--234; lines 288--294; lines 301--315; lines 554--558; lines 578--580; lines 602--612; lines 630--632.  Nothing was omitted from any retained span, so the package carries no omission record.  No retained line contains a citation, so the wrapper carries no bibliography and no bibliography entry.  No retained line contains a figure environment, an image-inclusion command or drawing code, so no image is delivered and the package declares no asset dependency.  Editorial formatting only, in addition: an AMS \texttt{amsart} preamble that restates the article's own declarations and macros used by the retained lines, namely the environments \texttt{lem}, \texttt{definition}, \texttt{thm} on separate counters, together with the standard packages those lines need.  The retained environments print in this document as Lemma 1, Definition 1, Theorem 1 in order of appearance, whereas the article prints the same items with the numbering of its own full document.  The complete native source of the article as distributed in the arXiv e-print for this exact version is bundled unchanged as \texttt{sources/171337/source0.tex}.
\begin{equation}\label{eq:Luf2}   
\Linf u=f  \textup{ in }  \Omega.
\end{equation}

\begin{lem}
\label{lem:strongmaxprinc}
   Let $u$ be a viscosity supersolution to $\Linf u = f$ in $\Omega$. Furthermore, assume that there exists some $x_0 \in \Omega$ such that $u(y) \geq u(x_0)$ for all $y \in \Rn \setminus \Omega$. Then $u$ is a constant function.
\end{lem}

\begin{lem}
\label{lem:Linfcont}
Let $u$ be a bounded viscosity supersolution to $\Linf u = f$ in $\Omega$. Then, for all $x \in \Omega$, we have that
\begin{equation}
\label{eq:Linf}
    \Lneg u(x) = \inf_{y \in \Rn \setminus \Omega} \frac{u(y)-u(x)}{|y-x|^{\alpha}},
\end{equation}
i.e., the infimum is attained in the complement of $\Omega$.
\end{lem}

\begin{definition}
\label{def:infconv}
    Let $u$ be bounded and lower semicontinuous. For $\epsilon > 0$ and $x \in \Rn$, we define the infimal convolution of $u$ as 
\begin{equation}\label{eq:infconv}
      u_{\epsilon}(x) \coloneqq \inf_{y \in \Rn} \left(\frac{|x-y|^2}{2 \epsilon} + u(y)\right) .    
\end{equation}
\end{definition}

\begin{lem}
\label{lem:infconv}
Let $u(x)$ be a viscosity supersolution to \eqref{eq:Luf2} such that $0 \leq u(x) \leq L$, for $x \in \Omega$. We take $\epsilon >0$ and define $\Omega_{\epsilon}$ as 
$$
\Omega_{\epsilon} = \{x \in \Omega: \dist(x, \partial \Omega)> r(\epsilon)\},
$$
where $r(\epsilon) \coloneqq \sqrt{2L\epsilon}$. Then, $u_{\epsilon}$ is a supersolution of
    \begin{equation*}
        \Linf \ue (x) = \tilde{f}^{\epsilon}(x),
    \end{equation*}
in $\Omega_{\epsilon}$ in the classical sense where
$$
\tilde{f}^{\epsilon}(x) =  f(x^*).
$$
\end{lem}

\begin{lem}
\label{lem:A1}
    Let $u(x)$ be bounded on $ \Rn$.
Then $M < \infty$ and $K_0 \subset \subset \Omega$.
\end{lem}

\begin{lem}\label{lem:A2}
      For some small enough $\beta>0$ we consider a compact $\beta$-neighbourhood $K(\beta)$ of $K_0$. Then $K_{\epsilon} \subset \subset K(\beta)$ for sufficiently small values of $\epsilon$.
\end{lem}

\begin{lem}\label{lem:A3}
    There exist an $\epsilon_0>0$ and a $\tau_0 >0$ such that for $\epsilon < \epsilon_0$, and $\tau \leq \tau_0$ we have that
\begin{equation}
\label{eq:blahblah2}
v \leq u_{\epsilon} + M_{\epsilon}/8 \ \textup{ on } \ \Rn \setminus \Omega_{\tau},   
\end{equation}
where
$$
\Omega_{\tau}:= \{ x \in \Omega, d(x, \partial \Omega)> \tau \}. 
$$
\end{lem}

\begin{lem}\label{lem:A4}
    Let $u(x)$ be bounded on $ \Rn$ and let the sequence $x_{\epsilon} $ be such that $x_{\epsilon} \rightarrow x_0$ for some $x_0 \in K_0$. Furthermore, assume that $\Lneg u_{\epsilon}(x_{\epsilon})$ is attained for some $y_{\epsilon} \in \Rn \setminus \Omega_{\epsilon}$ and that $y_{\epsilon} \to \infty$ as $\epsilon \to 0$. Then $u$ is a constant function.
\end{lem}

\begin{thm}\label{thm:comparison}
    Let $u$ and $v$ be a viscosity supersolution and a viscosity subsolution to \eqref{eq:Luf2}, respectively, where $0 < \alpha  \leq 1$, and $f \leq 0$ is bounded and continuous in $\Omega$. Furthermore, let $u(x)$ be bounded on $ \Rn \setminus \Omega$ and $u \geq v$ on $\Rn \setminus \Omega$. In addition, we assume that $u(x)$ and $v(x)$ have limits, as $x$ tends to infinity, in the sense that $\liminf_{|x| \to \infty} u(x) \geq C_1$ and $\limsup_{|x| \to \infty} v(x) \leq C_2$, where $C_1 \geq C_2$. Then $u \geq v$ in $\Omega$.
\end{thm}

\begin{proof}[Proof of Theorem \ref{thm:comparison}] We divide the proof into three steps.

\textbf{Step 1. Preliminaries} \\
At first we assume that $u$ is bounded in $\Omega$ i.e., $|u(x)| \leq \wt{L}$ for some constant $\wt{L} > 0$ whenever $x \in \Omega$. By adding a constant (to both $u$ and $v$) we may assume that $0\leq u(x)\leq L$ for some constant $L > 0$. We define
\begin{align*}
    M = \sup_{x \in \Rn} (v(x)-u(x)),
\end{align*}
and set
\begin{align*}
    K_0 = \{x \in \Rn:v(x)-u(x) = M \}.
\end{align*}
For the sake of contradiction, we assume that
\begin{equation}\label{eq:contradiction2}
    M \geq \gamma > 0.
\end{equation}
It follows from Lemma \ref{lem:A1} that $M$ must be finite and $K_0 \subset \subset \Omega$.
For some small enough $\beta>0$ we consider a compact $\beta$-neighbourhood $K(\beta)$ of $K_0$ such that
$$
K_0 \subset \subset K(\beta) \subset \subset \Omega,
$$
and for $x \in \Rn$ we define the infimal convolution of $u$ as in Definition \ref{def:infconv}. We recall that as $\epsilon \to 0^+$ it follows that $u_{\epsilon}(x) \to u(x)$ from below and by Lemma~\ref{lem:infconv}, $u_{\epsilon}$ is a supersolution to $\Linf u_{\epsilon} =  \tilde{f}^{\epsilon}$ in 
$$
\Omega_{\epsilon} = \{x \in \Omega:\dist(x, \partial \Omega) > r(\epsilon) \},
$$
where $r(\epsilon) = \sqrt{2L\epsilon}$, and 
$$
 \tilde{f}^{\epsilon}(x) =  f(x^*) \leq \sup_{y \in B_{r(\epsilon)}(x)} f(y).
$$
Here $x^*$ is the point where the infimum is attained for the infimal convolution of $u$ at the point $x$. We choose $\epsilon$ small enough so that $K_0 \subset \subset K(\beta) \subset \subset \Omega_{\epsilon}$ and proceed to define
\begin{align*}
    M_{\epsilon} &= \sup_{x \in \Rn}(v(x)-u_{\epsilon}(x)) , \\
    K_{\epsilon} &= \{x \in \Rn  :v(x)-u_{\epsilon}(x) = M_{\epsilon} \}.
\end{align*}
It follows from Lemma \ref{lem:A2} that $K_{\epsilon} \subset \subset K(\beta)$ for sufficiently small values of $\epsilon$.

We next consider $x_{\epsilon} \in K_{\epsilon}$. Since $\underbrace{v(x_{\epsilon})-u_{\epsilon}(x_{\epsilon})}_{M_{\epsilon}} \geq v(x)-u_{\epsilon}(x)$ for all $x \in \Rn$ it follows that
$$
v(x) \leq u_{\epsilon}(x)+M_{\epsilon},
$$
for all $x \in \Rn$.
We set
\begin{align*}
 \varphi_{\epsilon} =  u_{\epsilon}+M_{\epsilon},
\end{align*}
and thus
\begin{align}
     \begin{cases}
                 v \leq  \varphi_{\epsilon}  \ \textup{ in  } \Rn, \\
                 v = \varphi_{\epsilon} \ \textup{ in  }  K_{\epsilon}.
     \end{cases}
\end{align}
According to Lemma \ref{lem:A3} there exist an $\epsilon_0>0$ and a $\tau_0 >0$ such that for $\epsilon < \epsilon_0$, and $\tau \leq \tau_0$ we have that
\begin{equation}
\label{eq:blahblah}
v \leq u_{\epsilon} + M_{\epsilon}/8 \ \textup{ on } \ \Rn \setminus \Omega_{\tau},   
\end{equation}
where
$$
\Omega_{\tau}:= \{ x \in \Omega, d(x, \partial \Omega)> \tau \}. 
$$

 

We choose $\tau$ so small that $K(\beta) \subset \subset \Omega_{\tau}$ and take $\epsilon$ such that $r(\epsilon) < \tau$ so that $\partial \Omega_{\epsilon} \subset V$ where $V = \Rn \setminus \Omega_{\tau}$. Let $\eta_{\epsilon}(x) \in C^{\infty}(\Rn)$ be such that
\begin{align*}
     \begin{cases}
                   \eta_{\epsilon} = 1 & \textup{on  } \ \Rn \setminus \Omega_{\epsilon}, \\
                    0 \leq \eta_{\epsilon} \leq 1 & \textup{in  } \Omega_{\epsilon} \cap V, \\
                    \eta_{\epsilon} =0  & \textup{in  } \ \Omega_{\epsilon} \setminus V = \Omega_{\tau},
     \end{cases}
\end{align*}
and note that
$$
\Omega_{\epsilon} \cap V = \{x \in \Omega: r(\epsilon) < d(x, \partial \Omega) < \tau \}.
$$
Next, we define
\begin{equation*}
     \tilde{\varphi}_{\epsilon} = 
                \varphi_{\epsilon} -\frac{M_{\epsilon}}{4} \eta_{\epsilon},
\end{equation*}
and note that $\tilde{\varphi}_{\epsilon}= \varphi_{\epsilon}$ on $\Omega_{\epsilon} \setminus V$. Furthermore, we recall that 
$$
K_{\epsilon} \subset \subset K(\beta) \subset \subset \Omega_{\tau} = \Omega_{\epsilon} \setminus V,
$$
so $\tilde{\varphi}_{\epsilon} = \varphi_{\epsilon}$ in $K_{\epsilon}$. From equation \eqref{eq:blahblah} we know that
$$
v(x) \leq u_{\epsilon}(x) + M_{\epsilon}/8=\varphi_{\epsilon}(x)-\frac{7 M_{\epsilon}}{8} < \tilde{\varphi}_{\epsilon}(x),
$$
when $x \in V$. Therefore
\begin{align*}
     \begin{cases}
                   v \leq \tilde{\varphi}_{\epsilon} &\textup{in  } \Rn, \\
                   v= \tilde{\varphi}_{\epsilon} &\textup{in  } \ K_{\epsilon},
     \end{cases}
\end{align*}
and we conclude that $\tilde{\varphi}_{\epsilon}$ is a test function for $v$ at $x_{\epsilon} \in K_{\epsilon}$, i.e.,
\begin{equation}
\label{eq:contradiction}
\Linf \tilde{\varphi}_{\epsilon}(x_{\epsilon}) \geq f(x_{\epsilon}).    
\end{equation}
\textbf{Step 2. Obtaining a strict supersolution}\\
Next, we note that
\begin{align}
    \Lpos \tilde{\varphi}_{\epsilon}(x_{\epsilon}) &= \sup_{y \in \Rn} \frac{\tilde{\varphi}_{\epsilon}(y)-\tilde{\varphi}_{\epsilon}(x_{\epsilon})}{|y-x_{\epsilon}|^{\alpha}} = \sup_{y \in \Rn} \frac{\varphi_{\epsilon}(y)-\varphi_{\epsilon}(x_{\epsilon})-\eta_{\epsilon} (y)M_{\epsilon}/4}{|y-x_{\epsilon}|^{\alpha}} \nonumber \\
    & \leq \Lpos \varphi_{\epsilon}(x_{\epsilon}). \label{eq:Lplus}
\end{align}
 Furthermore, 
\begin{align}
    \Lneg \tilde{\varphi}_{\epsilon}(x_{\epsilon}) &= \Lneg \left(\varphi_{\epsilon}-\frac{M_{\epsilon}}{4} \eta_{\epsilon} \right)(x_{\epsilon}) \nonumber \\  
    &= \inf_{y \in \Rn} \frac{\varphi_{\epsilon}(y)-\frac{M_{\epsilon}}{4}\eta_{\epsilon} (y)-\varphi_{\epsilon}(x_{\epsilon})+\frac{M_{\epsilon}}{4}\eta_{\epsilon} (x_{\epsilon})}{|y-x_{\epsilon}|^{\alpha}}\nonumber   \\
    &=\inf_{y \in \Rn} \frac{\varphi_{\epsilon}(y)-\varphi_{\epsilon}(x_{\epsilon})-\frac{M_{\epsilon}}{4}\eta_{\epsilon} (y)}{|y-x_{\epsilon}|^{\alpha}} \nonumber \\
    &\leq  \inf_{y \in \Rn \setminus \Omega_{\epsilon}} \frac{\varphi_{\epsilon}(y)-\varphi_{\epsilon}(x_{\epsilon})-\frac{M_{\epsilon}}{4}}{|y-x_{\epsilon}|^{\alpha}}. \nonumber  %\label{eq:alignedinf}
\end{align}
We recall that $u_{\epsilon}(x)$ is a supersolution so by definition $\varphi_{\epsilon}(x)$ is a supersolution. If $\Lneg \varphi_{\epsilon}(x_{\epsilon})$ is attained for some $y_{\epsilon} \in \Rn \setminus \Omega_{\epsilon}$ it follows that
\begin{equation}\label{eq:wantthis}
   \Lneg \tilde{\varphi}_{\epsilon}(x_{\epsilon}) \leq \inf_{y \in \Rn \setminus \Omega_{\epsilon}} \frac{\varphi_{\epsilon}(y)-\varphi_{\epsilon}(x_{\epsilon})-\frac{M_{\epsilon}}{4}}{|y-x_{\epsilon}|^{\alpha}} \leq  \Lneg \varphi_{\epsilon}(x_{\epsilon})-\delta_{\epsilon},
\end{equation}
where $\delta_{\epsilon} = \frac{M_{\epsilon}}{4|y_{\epsilon}-x_{\epsilon}|^{\alpha}}$. Note that in the second inequality we used that
$$
\inf_{y \in \Rn} \frac{\varphi_{\epsilon}(y)-\varphi_{\epsilon}(x_{\epsilon})}{|y-x_{\epsilon}|^{\alpha}} = \inf_{y \in \Rn \setminus \Omega_{\epsilon}} \frac{\varphi_{\epsilon}(y)-\varphi_{\epsilon}(x_{\epsilon})}{|y-x_{\epsilon}|^{\alpha}},
$$
which is due to Lemma \ref{lem:Linfcont} since $u_{\epsilon}$ is a supersolution in $\Omega_{\epsilon}$. Next, we show that there exists some $\tilde{\delta}>0$ such that
\begin{equation}
\label{eq:delta_bound}
    \delta_{\epsilon} = \frac{M_{\epsilon}}{4|y_{\epsilon}-x_{\epsilon}|^{\alpha}} > \tilde{\delta},
\end{equation}
for $x_{\epsilon} \in K_{\epsilon}$ and very small values of $\epsilon$. If this is not true we either have that $|y_{\epsilon}|=\infty$ for all $\epsilon$ or there exists a sequence $x_{\epsilon}$ such that the corresponding sequence $y_{\epsilon} \to \infty$ as $\epsilon \to 0$. We start with the case $|y_{\epsilon}|=\infty$ for all $\epsilon$ less than some $\epsilon_0$. If this is true, then
$$
 \inf_{y \in \Rn \setminus \Omega_{\epsilon}} \frac{u_{\epsilon}(y)-u_{\epsilon}(x_{\epsilon})}{|y-x_{\epsilon}|^{\alpha}} =0.
$$
But this implies that $u_{\epsilon}(x_{\epsilon}) \leq u_{\epsilon}(y)$ for all $y \in \Rn \setminus \Omega_{\epsilon}$ and from Lemma \ref{lem:strongmaxprinc} it follows that $u_{\epsilon}$ is a constant function for all $\epsilon < \epsilon_0$.We move on to the second case and therefore assume that there exists a sequence $x_{\epsilon}$ such that the corresponding sequence $y_{\epsilon} \to \infty$ as $\epsilon \to 0$. According to Lemma \ref{lem:A4} this means that the function $u$ is constant which implies that the function $u_{\epsilon}$ is also constant. Then
\begin{equation*}
\inf_{y \in \Rn \setminus \Omega_{\epsilon}} \frac{\varphi_{\epsilon}(y) - \varphi_{\epsilon}(x_{\epsilon})-M_{\epsilon}/4}{|y-x_{\epsilon}|^{\alpha}} = \inf_{y \in \Rn \setminus \Omega_{\epsilon}} \frac{-M_{\epsilon}/4}{|y-x_{\epsilon}|^{\alpha}} = -\frac{M_{\epsilon}/4}{\inf_{y \in \Rn \setminus \Omega_{\epsilon}} |y-x_{\epsilon}|^{\alpha}}.
\end{equation*}
Therefore, in this case \eqref{eq:wantthis} holds with
$$
\delta_{\epsilon} = \frac{M_{\epsilon}/4}{\inf_{y \in \Rn \setminus \Omega_{\epsilon}} |y-x_{\epsilon}|^{\alpha}} \to  \frac{M/4}{\inf_{y \in \Rn \setminus \Omega} |y-x_0|^{\alpha}} >0,
$$
for all $\epsilon$ less than some $\epsilon_0$.

In conclusion, from \eqref{eq:Lplus} and \eqref{eq:wantthis} it follows that for $\epsilon$ small enough, we can find a $\delta > 0$ such that
\begin{align*}
    \Linf \tilde{\varphi}_{\epsilon}(x_{\epsilon}) &\leq \Linf \varphi_{\epsilon}(x_{\epsilon})- \delta  =\Linf u_{\epsilon}(x_{\epsilon}) - \delta  \label{eq:rhs_cont}
     \leq  \tilde{f}^{\epsilon}(x_{\epsilon})-\delta  \\
    & =f(x_{\epsilon}^*) -\delta  \leq \sup_{y \in B_{r(\epsilon)}(x_{\epsilon})} f(y)-\delta \leq f(x_{\epsilon})-\tilde{\delta}, \nonumber
\end{align*}
for some $\delta \geq \tilde{\delta}>0$. Note that we have used the continuity of the function $f$ in the last inequality. Therefore,
\begin{equation}\label{eq:contradiction3}
   \Linf \tilde{\varphi}_{\epsilon}(x_{\epsilon}) < f(x_{\epsilon}), 
\end{equation}
and thus, $\tilde{\varphi}$ is a strict supersolution.

\textbf{Step 3. Conclusion and removing the boundedness assumption}\\
We note that \eqref{eq:contradiction3} contradicts \eqref{eq:contradiction} which implies that \eqref{eq:contradiction2} is false. This proves the theorem under the assumption that $u$ is bounded in $\Omega$. 

We are now going to remove that assumption. Since $u$ is bounded on $\Rn \setminus \Omega$, i.e., $|u(x)|\leq C$ for $x \in \Rn \setminus \Omega$, it follows from Lemma \ref{lem:strongmaxprinc} that $u$ is bounded from below in $\Omega$. We next choose a constant $K$ so large that $K \geq C$ and 
\begin{equation}\label{eq:infext}
  \inf_{y \in \Rn \setminus \Omega} \frac{u(y)-K}{|y-x_0|^{\alpha}} \leq f(x_0),  
\end{equation}
for all $x_0 \in \Omega$. This can be done since $u$ is bounded on $\Rn \setminus \Omega$ and $f(x)$ is bounded. Indeed, we consider any $y \in \Rn \setminus \Omega$ and let $d_y$ be the maximum distance from $y$ to $\Omega$, i.e., $d_y = \sup_{x \in \Omega} |y-x|$. Then we can for instance choose $K$ such that
$$
\frac{u(y)-K}{d^{\alpha}_y} \leq \inf_{x \in \Omega} f(x).
$$
We next define
\begin{align*}
   \tilde{u}(x) =   \begin{cases}
                    \textup{min}\{u, K\} & \textup{in  } \ \Omega, \\
                    u    & \textup{on  } \ \Rn \setminus \Omega.
     \end{cases}
\end{align*}
The function $\tilde{u}$ is lower semicontinuous since $u$ is, the minimum of $u$ and $K$ is and $K \geq u$ on $\partial \Omega$. Furthermore, due to the choice of $K$ in \eqref{eq:infext}, $\tilde{u}$ is a viscosity supersolution to \eqref{eq:Luf2} with $f$ on the right hand side. Since it is bounded we can use the proof above to see that $v \leq \tilde{u}$ and therefore $v \leq u$. This concludes the proof of the theorem.
\end{proof}

\end{document}
